% GENERATED FILE. Edit canonical lessons, then run npm run generate:books.
% canonical-source-hash: fnv1a64:aff6bcf616fcea85
% canonical-lessons: MR-C354-nagarik, MR-C354-janta, MR-C354-sarvasamanya, MR-C354-loksankhya, MR-C354-gramin

\chapter{Citizens and the Public}
\label{ch:mr-citizens-and-the-public}
\hlchaptermodality{354}

\begin{chapteropening}
\textbf{By the end of this chapter:} \emph{I can talk about citizens, the public and rural life.}

Complete the last lesson of chapter 354: I can talk about citizens, the public and rural life.
\end{chapteropening}

\section[\texorpdfstring{\mr{नागरिक} (nāgarik)}{नागरिक (nāgarik)}]{\mr{नागरिक} (nāgarik) — a citizen}

\label{lesson:MR-C354-nagarik}



\begin{quote}
Before the new one: say the Marathi for actual, then the Marathi for to react.
\end{quote}

\subsection*{You'll want to know: \mr{नागरिक}}
\textbf{\mr{नागरिक}} — \emph{nāgarik} — \textquotedblleft{}a citizen\textquotedblright{}.

\begin{quote}
\mr{प्रत्येक} \mr{नागरिकाला} \mr{मत} \mr{देण्याचा} \mr{हक्क} \mr{आहे}.
\end{quote}

\begin{quote}
\emph{pratyek nāgarikālā mat deṇyācā hakka āhe.}
\end{quote}

\begin{quote}
Every citizen has the right to vote.
\end{quote}

Masculine.

\mr{नागरिक}, both a man and a woman; \mr{नागरिकत्व}, citizenship.

\begin{cousinweb}
\textbf{\mr{नागरिक}} (\emph{nāgarik}) is a tatsama, taken from Sanskrit \emph{nāgarika}, \textquotedblleft{}born or living in a town, civic\textquotedblright{} (Monier-Williams).
\end{cousinweb}

\begin{grammarlens}[title={What you've built}]
One more word for giving reasons and opinions.
\end{grammarlens}

\subsection*{Your turn}
\begin{itemize}
\raggedright
  \item \emph{Say it:} \emph{nāgarik}
  \item \emph{Say it:} \emph{nāgarik}, once more
  \item \emph{Say it:} \emph{pratyek nāgarikālā mat deṇyācā hakka āhe.} — Every citizen has the right to vote.
  \item \emph{Recall:} say the Marathi for actual, then the Marathi for to react, then say \emph{nāgarik} again
\end{itemize}

\begin{tcolorbox}[breakable,colback=teal!4,colframe=teal!35!black,title={Before you move on}]
What does \mr{नागरिक} mean? (\textquotedblleft{}A citizen\textquotedblright{}.) Say it once more.
\end{tcolorbox}
\section[\texorpdfstring{\mr{जनता} (jantā)}{जनता (jantā)}]{\mr{जनता} (jantā) — the general public; the people}

\label{lesson:MR-C354-janta}



\begin{quote}
Before the new one: say the Marathi for to react, then the Marathi for a citizen.
\end{quote}

\subsection*{You'll want to know: \mr{जनता}}
\textbf{\mr{जनता}} — \emph{jantā} — \textquotedblleft{}the general public; the people\textquotedblright{}.

\begin{quote}
\mr{नव्या} \mr{पुलाचा} \mr{जनतेला} \mr{खूप} \mr{फायदा} \mr{झाला}.
\end{quote}

\begin{quote}
\emph{navyā pulācā jantelā khūp phāydā jhālā.}
\end{quote}

\begin{quote}
The public benefited greatly from the new bridge.
\end{quote}

Feminine.

Feminine and singular, though it means many people: \mr{जनता} \mr{जमली}.

\begin{cousinweb}
\textbf{\mr{जनता}} (\emph{jantā}) is a tatsama, taken from Sanskrit \emph{janatā}, \textquotedblleft{}a number of people, the people\textquotedblright{} (Monier-Williams).
\end{cousinweb}

\begin{grammarlens}[title={What you've built}]
One more word for giving reasons and opinions.
\end{grammarlens}

\subsection*{Your turn}
\begin{itemize}
\raggedright
  \item \emph{Say it:} \emph{jantā}
  \item \emph{Say it:} \emph{jantā}, once more
  \item \emph{Say it:} \emph{navyā pulācā jantelā khūp phāydā jhālā.} — The public benefited greatly from the new bridge.
  \item \emph{Recall:} say the Marathi for to react, then the Marathi for a citizen, then say \emph{jantā} again
\end{itemize}

\begin{tcolorbox}[breakable,colback=teal!4,colframe=teal!35!black,title={Before you move on}]
What does \mr{जनता} mean? (\textquotedblleft{}The general public; the people\textquotedblright{}.) Say it once more.
\end{tcolorbox}
\section[\texorpdfstring{\mr{सर्वसामान्य} (sarvasāmānya)}{सर्वसामान्य (sarvasāmānya)}]{\mr{सर्वसामान्य} (sarvasāmānya) — common; ordinary}

\label{lesson:MR-C354-sarvasamanya}



\begin{quote}
Before the new one: say the Marathi for a citizen, then the Marathi for the general public.
\end{quote}

\subsection*{You'll want to know: \mr{सर्वसामान्य}}
\textbf{\mr{सर्वसामान्य}} — \emph{sarvasāmānya} — \textquotedblleft{}common; ordinary\textquotedblright{}.

\begin{quote}
\mr{सर्वसामान्य} \mr{माणसाला} \mr{हा} \mr{कायदा} \mr{समजत} \mr{नाही}.
\end{quote}

\begin{quote}
\emph{sarvasāmānya māṇsālā hā kāydā samjat nāhī.}
\end{quote}

\begin{quote}
The ordinary person does not understand this law.
\end{quote}

\mr{सर्व} + \mr{सामान्य}, all + common. Invariable.

\begin{cousinweb}
\textbf{\mr{सर्वसामान्य}} (\emph{sarvasāmānya}) is a tatsama, taken from Sanskrit \emph{sarvasāmānya}, \textquotedblleft{}common to all\textquotedblright{} (Monier-Williams).
\end{cousinweb}

\begin{grammarlens}[title={What you've built}]
One more word for giving reasons and opinions.
\end{grammarlens}

\subsection*{Your turn}
\begin{itemize}
\raggedright
  \item \emph{Say it:} \emph{sarvasāmānya}
  \item \emph{Say it:} \emph{sarvasāmānya}, once more
  \item \emph{Say it:} \emph{sarvasāmānya māṇsālā hā kāydā samjat nāhī.} — The ordinary person does not understand this law.
  \item \emph{Recall:} say the Marathi for a citizen, then the Marathi for the general public, then say \emph{sarvasāmānya} again
\end{itemize}

\begin{tcolorbox}[breakable,colback=teal!4,colframe=teal!35!black,title={Before you move on}]
What does \mr{सर्वसामान्य} mean? (\textquotedblleft{}Common; ordinary\textquotedblright{}.) Say it once more.
\end{tcolorbox}
\section[\texorpdfstring{\mr{लोकसंख्या} (loksaṅkhyā)}{लोकसंख्या (loksaṅkhyā)}]{\mr{लोकसंख्या} (loksaṅkhyā) — population}

\label{lesson:MR-C354-loksankhya}



\begin{quote}
Before the new one: say the Marathi for the general public, then the Marathi for common.
\end{quote}

\subsection*{You'll want to know: \mr{लोकसंख्या}}
\textbf{\mr{लोकसंख्या}} — \emph{loksaṅkhyā} — \textquotedblleft{}population\textquotedblright{}.

\begin{quote}
\mr{पुण्याची} \mr{लोकसंख्या} \mr{दरवर्षी} \mr{वाढत} \mr{आहे}.
\end{quote}

\begin{quote}
\emph{puṇyācī loksaṅkhyā darvarṣī vāḍhat āhe.}
\end{quote}

\begin{quote}
Pune's population is growing every year.
\end{quote}

Feminine.

\mr{लोक} + \mr{संख्या}, people + number.

\begin{cousinweb}
\textbf{\mr{लोकसंख्या}} (\emph{loksaṅkhyā}) is a learned compound built from Sanskrit \emph{loka}, \textquotedblleft{}world, people\textquotedblright{}, and \emph{saṁkhyā}, \textquotedblleft{}number\textquotedblright{} (Monier-Williams).
\end{cousinweb}

\begin{grammarlens}[title={What you've built}]
One more word for giving reasons and opinions.
\end{grammarlens}

\subsection*{Your turn}
\begin{itemize}
\raggedright
  \item \emph{Say it:} \emph{loksaṅkhyā}
  \item \emph{Say it:} \emph{loksaṅkhyā}, once more
  \item \emph{Say it:} \emph{puṇyācī loksaṅkhyā darvarṣī vāḍhat āhe.} — Pune's population is growing every year.
  \item \emph{Recall:} say the Marathi for the general public, then the Marathi for common, then say \emph{loksaṅkhyā} again
\end{itemize}

\begin{tcolorbox}[breakable,colback=teal!4,colframe=teal!35!black,title={Before you move on}]
What does \mr{लोकसंख्या} mean? (\textquotedblleft{}Population\textquotedblright{}.) Say it once more.
\end{tcolorbox}
\section[\texorpdfstring{\mr{ग्रामीण} (grāmīṇ)}{ग्रामीण (grāmīṇ)}]{\mr{ग्रामीण} (grāmīṇ) — rural}

\label{lesson:MR-C354-gramin}



\begin{quote}
Before the new one: say the Marathi for common, then the Marathi for population.
\end{quote}

\subsection*{You'll want to know: \mr{ग्रामीण}}
\textbf{\mr{ग्रामीण}} — \emph{grāmīṇ} — \textquotedblleft{}rural\textquotedblright{}.

\begin{quote}
\mr{ग्रामीण} \mr{भागात} \mr{अजून} \mr{चांगले} \mr{रस्ते} \mr{नाहीत}.
\end{quote}

\begin{quote}
\emph{grāmīṇ bhāgāt ajūn cāṅgle raste nāhīt.}
\end{quote}

\begin{quote}
There are still no good roads in rural areas.
\end{quote}

From \mr{ग्राम}, a village (formal); the opposite is \mr{शहरी}, urban.

\begin{cousinweb}
\textbf{\mr{ग्रामीण}} (\emph{grāmīṇ}) is a tatsama, taken from Sanskrit \emph{grāmīṇa}, \textquotedblleft{}rural, of a village\textquotedblright{} (Monier-Williams).
\end{cousinweb}

\begin{grammarlens}[title={What you've built}]
One more word for giving reasons and opinions.
\end{grammarlens}

\subsection*{Your turn}
\begin{itemize}
\raggedright
  \item \emph{Say it:} \emph{grāmīṇ}
  \item \emph{Say it:} \emph{grāmīṇ}, once more
  \item \emph{Say it:} \emph{grāmīṇ bhāgāt ajūn cāṅgle raste nāhīt.} — There are still no good roads in rural areas.
  \item \emph{Recall:} say the Marathi for common, then the Marathi for population, then say \emph{grāmīṇ} again
\end{itemize}

\begin{tcolorbox}[breakable,colback=teal!4,colframe=teal!35!black,title={Before you move on}]
What does \mr{ग्रामीण} mean? (\textquotedblleft{}Rural\textquotedblright{}.) Say it once more.
\end{tcolorbox}
